\documentclass[12pt]{article}
% Préambule commun à tous les documents extraits de « La Revue CAPES »
\usepackage[T1]{fontenc}
\usepackage[utf8]{inputenc}
\usepackage{lmodern}
\usepackage{amsmath,amssymb,amsthm,amsfonts,mathrsfs}
\usepackage{setspace}
\usepackage{listings}
\usepackage{pgfplots}
\pgfplotsset{compat=1.18}
\usepackage{longtable}
\usepackage{adjustbox}
\usepackage{lettrine}
\usepackage{tkz-tab}
\usepackage{changepage}
\usepackage{pdflscape}
\usepackage{graphicx}
\usepackage{tikz}
\usepackage{xcolor}
\usepackage[a4paper,top=2cm,bottom=2.5cm,left=2.5cm,right=2cm]{geometry}
\usepackage{fancyhdr}
\usepackage{draftwatermark}
\SetWatermarkText{La RevueCapes}
\SetWatermarkLightness{0.9}
\SetWatermarkScale{2}
\usepackage{hyperref}
\hypersetup{colorlinks=true,linkcolor=blue!50!black,urlcolor=blue!50!black}

\definecolor{revue}{RGB}{20,60,120}

\pagestyle{fancy}
\fancyhf{}
\renewcommand{\headrulewidth}{0.4pt}
\setlength{\headheight}{15pt}

% \entete{Matière}{Titre}{Type de document}
\newcommand{\entete}[3]{%
  \fancyhead[L]{\small\textsc{La Revue CAPES} -- #1}%
  \fancyhead[R]{\small #2}%
  \fancyfoot[C]{\thepage}%
  \thispagestyle{plain}%
  \begin{center}
    {\color{revue}\rule{\textwidth}{1.2pt}}\\[4pt]
    {\small\textsc{Concours directs CAP-CEG / CAPES -- Burkina Faso}}\\[6pt]
    {\LARGE\bfseries\color{revue} #2}\\[6pt]
    {\large\itshape #1 -- #3}\\[2pt]
    {\color{revue}\rule{\textwidth}{1.2pt}}
  \end{center}
  \vspace{0.5cm}%
}

% Encadré pour les remarques ajoutées dans les corrigés
\newcommand{\remarque}[1]{\par\medskip\noindent\fcolorbox{revue}{revue!6}{\parbox{\dimexpr\linewidth-2\fboxsep-2\fboxrule}{\textbf{\color{revue}Remarque.} #1}}\par\medskip}


\begin{document}
\entete{Mathématiques}{Sujet type n°6}{Énoncé}
\begin{spacing}{1.5}

 
 \medskip\textbf{Exercice 1 }

 Pour tout $x,y \in \mathbb{R}$, montrer les inégalités suivantes :
 
 1. Première inégalité triangulaire $\vert x+y\vert \leq \vert x\vert + \vert y\vert$
 
 2. Seconde inégalité triangulaire $\vert\vert x\vert - \vert y\vert\vert\leq\vert x-y\vert$.
 
 \medskip\textbf{Exercice 2}

 Soient $X$ et $Y$ deux variables aléatoires dont les lois sont données par :
 
 $P(X =1)=0.6$, $P(Y =1)=0.4$
 
 $P(X =2)=0.2$, $P(Y =2)=0.5$
 
$P(X =3)=0.2$, $P(Y =3)=0.1$
 
 1. La loi du couple $(X,Y)$ est donnée par l'un des tableaux suivants. Lequel?
 
 \begin{center}
\begin{tabular}{c}\textbf{Tableau A}\\[2pt]
\begin{tabular}{|c|c|c|c|}
 \hline 
 $Y\setminus X $& 1 & 2 & 3 \\ 
 \hline 
 1 & 0,5 & 0 & 0,1 \\ 
 \hline 
 2 & 0,2 & 0,1 & 0 \\ 
 \hline 
 3 & 0,1 & 0 & 0 \\ 
 \hline 
 \end{tabular}
\end{tabular}\qquad
\begin{tabular}{c}\textbf{Tableau B}\\[2pt]
\begin{tabular}{|c|c|c|c|}
 \hline 
 $Y\setminus X $& 1 & 2 & 3 \\ 
 \hline 
 1 & 0,5 & 0 & 0,1 \\ 
 \hline 
 2 & 0 & 0,1 & 0 \\ 
 \hline 
 3 & 0,1 & 0 & 0 \\ 
 \hline 
 \end{tabular}
\end{tabular}\\[10pt]
\begin{tabular}{c}\textbf{Tableau C}\\[2pt]
\begin{tabular}{|c|c|c|c|}
 \hline 
 $Y\setminus X $& 1 & 2 & 3 \\ 
 \hline 
 1 & 0,2 & 0,1 & 0,1 \\ 
 \hline 
 2 & 0,3 & 0,1 & 0,1 \\ 
 \hline 
 3 & 0,1 & 0 & 0 \\ 
 \hline 
 \end{tabular}
\end{tabular}\qquad
\begin{tabular}{c}\textbf{Tableau D}\\[2pt]
\begin{tabular}{|c|c|c|c|}
 \hline 
 $Y\setminus X $& 1 & 2 & 3 \\ 
 \hline 
 1 & 1 & 0 & 0 \\ 
 \hline 
 2 & 0 & 1 & 0 \\ 
 \hline 
 3 & 0 & 0 & 1 \\ 
 \hline 
 \end{tabular}
\end{tabular}\\[10pt]
\end{center}

2. Quel serait le tableau de la loi du couple $(X,Y)$ si les variables $X$ et $Y$ étaient
 indépendantes ?

3. Parmi les autres tableaux, déterminer lesquels peuvent correspondre à la loi d'un couple de variables. Lorsque c'est le cas, donner les lois marginales correspondantes.

\medskip\textbf{Exercice 3}

On considère les deux fonctions $f_n$ et $g_n$, définies sur $\mathbb{R}$, par : 
$f_n(x) = (1 + x)^n$ 

$g_n(x) = (1 - x)^n$ où $n$ est un entier naturel. 

1) Par un calcul direct, donner les expressions de leurs primitives, notées respectivement $F_n(x)$ 
et $G_n(x)$.  

2) En développant $f_n$ et $g_n$ par la formule du binôme, donner deux autres expressions de $F_n(x)$ et $G_n(x)$. 

3) Calculer les quantités $F_n(1)-F_n(0)$ et $G_n(1) - G_n(0)$. 

En déduire la valeur de $A$ et $B$ où : 

\[
A = \sum_{k=0}^{k=n}\frac{1}{k+1}C_{n}^{k}
\]

\[
B=\sum_{k=0}^{k=n}(-1)^{k+1}\frac{1}{k+1}C_{n}^{k}
\]

\medskip\textbf{Exercice 4}

 On munit $\mathbb{R}^n$ de son produit scalaire usuel.
 
 1. Montrer que pour tout $x = (x_i) \in \mathbb{R}^n$, on a
 \[
 (\sum_{i=1}^{n} x_i)^2\leq n(\sum_{i=1}^{n} x_{i}^{2})
 \]
 
 2.  Soit $A \in M_{n}(\mathbb{R}), A = [a_{ij}]$.
 
 Calculer $tr(A^tA)$ en fonction des $a_{ij}$.
 
 En déduire l'inégalité $\mid tr(A)\mid \leq \sqrt{ntr(A^tA)}$ pour toute matrice $A$.
 
 3. Montrer que sur $M_n(\mathbb{R})$ l'application $A,B \longmapsto tr(A^tB)$ définit un produit scalaire.
 
4. Montrer que pour toutes matrices $A$ et $B$ symétriques, on a :
 
 $\mid tr(AB)\mid^2 \leq  (tr(A^2))(tr(B^2))$.
 
 \quad
 

\end{spacing}
\end{document}
